\documentclass[../../main.tex]{subfiles}
\begin{document}

\paragraph{【新課程】}
2つの関数(函数) $y_1=\dfrac{\pi}{4}\dfrac{1}{x}-\dfrac{\pi}{4}+1$ と
$y_2=\sqrt{2}\cos\dfrac{\pi}{4}x$ とのグラフを考える．
\begin{enumerate}
  \item この2つのグラフは，$0<x \leqq 1$の範囲では，
    点$(1,1)$以外に交点がないことを次の方針で示せ．\\
    $y_3=\dfrac{\pi}{4}(1-x)+1$を$0<x \leqq 1$の範囲で考えると，
    $y_1 \geqq y_3 \geqq y_2$であり，等号は$x=1$のときに限り成立する．
  \item 範囲$0 \leqq x \leqq 1$において，
    $y_1$のグラフ，$y_2$のグラフ，$y$軸，直線$y=\dfrac{\pi}{4}+1$
    で囲まれた部分の面積を求めよ．
\end{enumerate}

\paragraph{【旧課程】}
長さの単位をセンチメートル，時間$t$の単位を秒とする．
曲線$y=x^2$の軸を鉛直にして，この曲線を軸のまわりに回転してえられる曲面を内面とする容器がある．
ある時刻$(t=0)$に水をこの容器に入れ始め，任意の$t$ $(>0)$に対して，
$t$秒後の水面の上昇速度が$t^2\,\text{cm/sec}$であるようにするには，
水の注入速度(単位は$\text{cm}^3/\text{sec}$)をどのようにすればよいか．

\end{document}
